\documentclass[11pt,a4paper]{article}
\usepackage[utf8]{inputenc}
\usepackage{amsmath,amssymb,amsthm}
\usepackage{mathtools}
\usepackage{geometry}
\geometry{margin=2.5cm}

\newtheorem{theorem}{Theorem}[section]
\newtheorem{lemma}[theorem]{Lemma}
\newtheorem{proposition}[theorem]{Proposition}
\newtheorem{corollary}[theorem]{Corollary}
\theoremstyle{definition}
\newtheorem{definition}[theorem]{Definition}
\theoremstyle{remark}
\newtheorem{remark}[theorem]{Remark}

\newcommand{\N}{\mathbb{N}}
\newcommand{\R}{\mathbb{R}}
\newcommand{\C}{\mathbb{C}}
\newcommand{\Z}{\mathbb{Z}}
\newcommand{\abs}[1]{\lvert#1\rvert}
\newcommand{\norm}[1]{\lVert#1\rVert}
\newcommand{\set}[1]{\{#1\}}
\newcommand{\p}{\operatorname{p}}
\newcommand{\GDST}{\operatorname{GDST}}

\title{Proof of Part III: The Greedy Dirichlet Sub‑Sum Transform}
\author{Based on the formalisation in \texttt{GDST\_Transform.thy}}
\date{}

\begin{document}
\maketitle

\begin{abstract}
We define the Greedy Dirichlet Sub‑Sum Transform (GDST) \(E_{\GDST}(\theta,s)\), which is the Dirichlet series formed by restricting the harmonic expansion of \(\theta\) to those indices where the greedy digit equals \(1\).
We prove that for any \(\theta\in[0,\pi]\) the series converges absolutely on the half‑plane \(\{s\in\C:\Re s>0\}\) and defines a holomorphic function there.
The complement series \(\Omega_{\GDST}(\theta,s)\) collects the remaining terms, and together they satisfy
\[
\zeta(s)-1 = E_{\GDST}(\theta,s) + \Omega_{\GDST}(\theta,s) \qquad (\Re s>1),
\]
where \(\zeta\) is the Riemann zeta function.
\end{abstract}

\section{The GDST transform: definition and holomorphy}\label{sec:GDST_def}

Let \(\theta\in[0,\pi]\) and let \(x=\theta/\pi\in[0,1]\) be its normalised coordinate.
Recall from Part~I the greedy digits \(\delta_n(\theta) = \delta_n(x)\) and the greedy expansion
\[
x = \sum_{n=0}^{\infty}\frac{\delta_n(x)}{n+2}.
\]

\begin{definition}[GDST transform]\label{def:E_gdst}
For \(\theta\in[0,\pi]\) and \(s\in\C\) with \(\Re s>0\) define
\begin{equation}\label{eq:E_def}
E_{\GDST}(\theta,s) = \sum_{n=0}^{\infty} \delta_n(\theta)\,(n+2)^{-s}.
\end{equation}
\end{definition}

In the series, only those indices \(n\) with \(\delta_n(\theta)=1\) contribute; the convergence is absolute and uniform on compact subsets of \(\{\Re s>0\}\), as we shall prove.

\subsection{Absolute convergence}

The absolute convergence of \(E_{\GDST}\) rests on the super‑exponential growth of selected indices (Theorem~selected\_growth).

\begin{lemma}[summability of the GDSt terms]\label{lem:summ_GDST}
For any \(\theta\in[0,\pi]\) and any \(s\in\C\) with \(\Re s>0\), the series
\[
\sum_{n=0}^{\infty} \delta_n(\theta)\,(n+2)^{-s}
\]
converges absolutely.
\end{lemma}
\begin{proof}
Write \(x=\theta/\pi\). By Part~I, Corollary~summable\_delta\_powr, the series
\[
\sum_{n=0}^{\infty} \delta_n(x)\,(n+2)^{-\Re s}
\]
converges (the corollary was proved for all \(x\in[0,1]\) and \(\Re s>0\)). 
Since \(\delta_n(\theta)=\delta_n(x)\) and \(\abs{(n+2)^{-s}} = (n+2)^{-\Re s}\), we have
\[
\abs{\delta_n(\theta)\,(n+2)^{-s}} = \delta_n(x)\,(n+2)^{-\Re s},
\]
the same terms as in the real series. Hence absolutely summability follows immediately.
\end{proof}

\subsection{Holomorphy}

\begin{theorem}[holomorphy of \(E_{\GDST}\)]\label{thm:holo}
For fixed \(\theta\in[0,\pi]\), the function \(s\mapsto E_{\GDST}(\theta,s)\) is holomorphic on the open half‑plane \(\{s\in\C : \Re s>0\}\).
\end{theorem}
\begin{proof}
We use Weierstrass’ theorem on uniform convergence of series of holomorphic functions.
For each \(N\in\N\), the partial sum
\[
F_N(s) = \sum_{n=0}^{N} \delta_n(\theta)\,(n+2)^{-s}
\]
is a finite linear combination of exponentials, hence entire (holomorphic on \(\C\)).

Let \(K\subset\{s\mid\Re s>0\}\) be compact. Set
\[
\delta = \inf_{s\in K}\Re s \;>\;0.
\]
Then for every \(s\in K\) and every \(n\ge0\),
\[
\abs{\delta_n(\theta)\,(n+2)^{-s}} = \delta_n(\theta)\,(n+2)^{-\Re s} \le \delta_n(\theta)\,(n+2)^{-\delta}.
\]
By Lemma~\ref{lem:summ_GDST} (with \(s=\delta\)), the series of non‑negative terms
\[
\sum_{n=0}^{\infty} \delta_n(\theta)\,(n+2)^{-\delta}
\]
is convergent. Consequently,
\[
M_n := \delta_n(\theta)\,(n+2)^{-\delta}
\]
is a summable sequence that dominates \(\abs{\delta_n(\theta)\,(n+2)^{-s}}\) uniformly on \(K\).

By the Weierstrass \(M\)-test, the series \(\sum_{n=0}^{\infty} \delta_n(\theta)\,(n+2)^{-s}\) converges uniformly on \(K\). Since each partial sum is holomorphic and uniform limits of holomorphic functions on compact sets are holomorphic, the limit \(E_{\GDST}(\theta,\cdot)\) is holomorphic on the half‑plane.
\end{proof}

\section{The complement series and relation with \(\zeta\)}\label{sec:complement}

\begin{definition}[complement GDSt]\label{def:Omega}
For \(\theta\in[0,\pi]\) and \(\Re s>0\) define the complement series
\begin{equation}\label{eq:Omega_def}
\Omega_{\GDST}(\theta,s) = \sum_{n=0}^{\infty} \bigl(1-\delta_n(\theta)\bigr)\,(n+2)^{-s}.
\end{equation}
\end{definition}
That is, \(\Omega_{\GDST}\) sums the terms with \(\delta_n=0\). By exactly the same argument as for \(E_{\GDST}\), this series converges absolutely and defines a holomorphic function on \(\{\Re s>0\}\).

\begin{theorem}[zeta‑split]\label{thm:zeta_split}
For every \(\theta\in[0,\pi]\) and every \(s\in\C\) with \(\Re s>1\),
\[
\zeta(s) - 1 = E_{\GDST}(\theta,s) + \Omega_{\GDST}(\theta,s).
\]
\end{theorem}
\begin{proof}
The Riemann zeta function can be written (for \(\Re s>1\)) as
\[
\zeta(s) = \sum_{n=1}^{\infty} \frac{1}{n^s} = 1 + \sum_{n=2}^{\infty} \frac{1}{n^s}.
\]
Re‑indexing with \(n\mapsto n+2\) gives
\[
\zeta(s)-1 = \sum_{n=0}^{\infty} \frac{1}{(n+2)^s}.
\]

For each integer \(n\ge0\), precisely one of the digits \(\delta_n(\theta)\) and \(1-\delta_n(\theta)\) equals \(1\), because \(\delta_n(\theta)\in\{0,1\}\). Hence
\[
\frac{1}{(n+2)^s} = \delta_n(\theta)\,(n+2)^{-s} + \bigl(1-\delta_n(\theta)\bigr)\,(n+2)^{-s}.
\]
Summing over all \(n\) and using the absolute convergence (which allows termwise addition) we obtain
\[
\zeta(s)-1 = \sum_{n=0}^{\infty} \delta_n(\theta)\,(n+2)^{-s} + \sum_{n=0}^{\infty} \bigl(1-\delta_n(\theta)\bigr)\,(n+2)^{-s} = E_{\GDST}(\theta,s) + \Omega_{\GDST}(\theta,s),
\]
as claimed.
\end{proof}

\begin{remark}
The identity shows that \(E_{\GDST}\) and \(\Omega_{\GDST}\) provide a partition of the tail of the Riemann zeta function, indexed by the greedy digits. Because the selected indices grow super‑exponentially, \(E_{\GDST}\) is a sparse subseries of the full zeta series, and its zero structure is intimately tied to the zeros of \(\zeta\) — a connection that will be exploited via the transfer operator in the following parts.
\end{remark}

\end{document}